\subsection{What continual training changes}
\label{sec:bf-is-craft-detail}

An accurate transport between two neighboring temperatures need not be an
accurate transport from the prior to the final posterior. This observation
motivates learning separate maps for the easier intermediate problems.
For maps $T_{k,\phi_k}$, CRAFT's population objective is
\begin{equation}
 \mathcal H(\phi_1,\ldots,\phi_K)
 =\sum_{k=1}^K D_{\rm KL}
       ((T_{k,\phi_k})_\#\pi_{k-1}\Vert\pi_k).
 \label{eq:bf-is-craft-objective}
\end{equation}
Here $T_\#\pi$ denotes the law of $T(X)$ when $X\sim\pi$.
Each term vanishes when its map transports the preceding distribution
exactly to the next one. This is a local reverse-KL objective, rather than
forward KL against a fixed final-posterior dataset or FAB's second-moment
objective \citep[Section 2.3, Equations 5--8]{matthews2022craft}.

To reconstruct the correction, let $y=T_k(x)$ and write
$q_k(y)=\pi_{k-1}(x)/|\det DT_k(x)|$. Then
\begin{equation}
 G_k(x)q_k(T_k(x))
 =\frac{\gamma_k(T_k(x))}{Z_{k-1}}
 =\frac{Z_k}{Z_{k-1}}\pi_k(T_k(x)).
 \label{eq:bf-is-craft-correction}
\end{equation}
Thus the factor in \eqref{eq:bf-is-aft-weight} corrects imperfect transport;
the unknown normalizers cancel from relative particle weights. An invariant
MCMC kernel applied afterward preserves the resulting target. With an
identity map the construction reduces to ordinary annealed SMC. With an
exact map every incremental weight is $Z_k/Z_{k-1}$, although subsequent
resampling can still create dependence among particles.

AFT learns a map from the finite cloud available at a particular temperature.
Repeated optimization against that cloud can fit its sampling errors.
Practical AFT uses separate training, validation and test populations to
limit overfitting, but a single pass still supplies limited new information
at each stage. CRAFT instead makes repeated passes from the base distribution.
During a pass it evaluates each map's loss and transports particles with
the current map. The newly computed optimizer update affects later use of
that map. The public implementation collects the stage gradients through
the pass and then updates the parameters. This ordering avoids fitting a
map to the very particles whose fixed-map importance identity is being
invoked \citep[Algorithms 1--3 and Section 2.3]{matthews2022craft}.

The distinction between a particle approximation and the ideal distribution
also matters for the gradient. Set
\begin{align}
 D_k(x;\phi_k)&=\log\gamma_{k-1}(x)
       -\log\gamma_k(T_{k,\phi_k}(x))
       -\log|\det DT_{k,\phi_k}(x)|,\nonumber\\
 \widehat g_k&=\sum_i W_{k-1,i}\nabla_{\phi_k}D_k(X_{k-1,i};\phi_k).
 \label{eq:bf-is-craft-local-gradient}
\end{align}
For fixed upstream parameters, define
$\widehat\pi_{k-1}^N=\E[\sum_i W_{k-1,i}\delta_{X_{k-1,i}}]$.
Under the differentiation and integrability conditions of
\citet[Proposition 2]{matthews2022craft}, $\widehat g_k$ is unbiased for
the local gradient of
$D_{\rm KL}((T_{k,\phi_k})_\#\widehat\pi_{k-1}^N\Vert\pi_k)$.
At finite $N$, $\widehat\pi_{k-1}^N$ generally differs from $\pi_{k-1}$.
The proposition therefore does not turn a self-normalized particle
gradient into an exactly unbiased gradient of the ideal population objective.
Nor should the local update be confused with differentiating every later
particle distribution through every earlier optimizer update.

The examples in \citet[Sections 4.1--4.3]{matthews2022craft} include a
1024-dimensional log-Gaussian Cox process, a 30-dimensional variational
autoencoder latent distribution, and a lattice-field application with a
particle-MCMC correction. Their VAE experiment uses affine inverse
autoregressive flows. These examples establish that the construction is
not restricted to tiny target dimension; they do not establish coverage
of a new posterior with different modes or recurrent-state dynamics.
For BayesFilter, the attraction is the combination of refreshed training
samples, short transport tasks and explicit importance correction. Its
remaining empirical question is whether those mechanisms discover and
retain the posterior regions that the current training misses.

Both the correction and the local loss require a forward map and its log
Jacobian, without evaluating an inverse-flow density at every transported
point. This is computationally relevant for IAF, whose forward direction
is parallel across coordinates while inversion is autoregressive. Target
values and training gradients still have their own cost. Moreover, the
complete sampler includes resampling and mutation: composing its stage
maps alone does not reproduce its output distribution or supply that
distribution's tractable density. A separate final IAF fitted to weighted
sampler output is therefore a distinct learned approximation.

\section{Adaptive mutation and waste-free SMC}
\label{sec:bf-is-modern-smc-mutation}

An annealing schedule controls changes in the target distribution. Mutation
controls how far the particle population explores at each temperature.
These are separate decisions: smaller temperature increments can preserve
balanced weights while particles remain confined to their original basins.
The posterior path $r(\theta)L(\theta)^\beta$ is already explicit in
\citet[Section 4.2.2]{delmoral2006smc}; modern adaptive methods improve how
effort is allocated along such a path.

\subsection{Tune movement relative to its cost}

For an HMC proposal $Y_h$ from $X$, let $h=(\epsilon,L)$ contain step size
and leapfrog count, and let $\alpha_h$ be its Metropolis acceptance
probability. With a declared positive-definite scaling matrix $A$, a useful
local tuning quantity is
\begin{equation}
 U(h)=\frac{\E[\alpha_h(X,Y_h)
                (Y_h-X)^\top A^{-1}(Y_h-X)]}{c(h)}.
 \label{eq:bf-is-smc-movement-cost}
\end{equation}
\citet[Section 3.2, Equation 6]{buchholz2021hmcsmc} use acceptance-weighted
squared jumping distance divided by $L$, treating leapfrog count as a proxy
for computational cost. They also investigate a pretuning procedure at each
temperature. The population supplies information about which step sizes and
trajectory lengths produce useful movement.

This criterion explains why acceptance alone is inadequate: arbitrarily small
proposals can be accepted almost always while moving almost nowhere.
It also explains why an expensive trajectory should not be rewarded solely
for a larger displacement. In a library implementation $c(h)$ should reflect
measured value/score work when the simple proportional-to-$L$ model is poor.
The scaling matrix determines which movements are counted as large, so its
choice and estimation are part of the method, not innocuous bookkeeping.

Neither $U(h)$ nor high acceptance measures communication between separated
modes. A chain can make large moves within one broad mode and still never
reach another. Adaptive kernel choices also require their own SMC
justification; the fixed-kernel identities above cannot be applied without
qualification to arbitrary population-dependent tuning. A practical way to
obtain interpretable evidence is to calibrate a ladder and kernels in pilot
populations, freeze them, and assess independent subsequent populations.
This is a study design, not a claim that all mathematically justified
adaptive SMC algorithms must freeze their schedules.

\subsection{Retain the intermediate mutation states}

Ordinary resample--move SMC may run a chain for several steps and retain
only its last state. Waste-free SMC resamples $M$ ancestors, runs a chain of
length $P$ from each, and retains all $N=MP$ states
\citep[Section 2.4, Algorithm 2]{dau2022wastefree}. In the paper's indexing,
the chain at stage $k$ preserves the preceding target before the next
potential is applied. With those conventions, its weighted estimate is
\begin{equation}
 \widehat\pi_k(f)=
 \frac{\sum_{m=1}^M\sum_{p=1}^P
                  G_k(X_k^{m,p})f(X_k^{m,p})}
      {\sum_{m=1}^M\sum_{p=1}^P G_k(X_k^{m,p})}.
 \label{eq:bf-is-wastefree-estimate}
\end{equation}
One cannot reproduce this algorithm merely by appending unweighted warm-up
states to an ordinary SMC output: ancestor selection, invariant target,
potential and ordering must match its Feynman--Kac construction.

The $MP$ states are not $MP$ independent observations. To illustrate the
effect of dependence, suppose temporarily that the $M$ chains are independent
and stationary for one distribution and that $f$ has summable
autocorrelations $\rho_f(\ell)$. For long chains,
\begin{equation}
 \Var\!\left[\frac{1}{MP}\sum_{m,p}f(X^{m,p})\right]
 \simeq\frac{\Var_\pi(f)}{MP}
       \left(1+2\sum_{\ell=1}^{\infty}\rho_f(\ell)\right).
 \label{eq:bf-is-wastefree-autocorrelation}
\end{equation}
This calculation explains why retaining intermediate states can use
computation more efficiently while leaving a substantial correlation
penalty. It is not the variance formula for the entire interacting SMC
history. That history needs the paper's asymptotic analysis or uncertainty
estimated from independent full populations.

The long-chain results in \citet[Section 3.1]{dau2022wastefree} assume
bounded potentials and uniformly ergodic mutation kernels. Their
identity-kernel example is instructive: all $P$ descendants of an ancestor
are identical, so increasing $P$ cannot eliminate the error of a fixed
$M$-point population. Their adaptive chain-length strategy estimates
autocorrelation and increases effort where mixing deteriorates. It reduces
wasted mutation effort; it does not turn a missing-mode problem into an
ordinary ESS calculation.

\section{Dimension, overlap and mixing}
\label{sec:bf-is-modern-smc-dimension}

Consider a product target $p_d(x)=\prod_{j=1}^d p_1(x_j)$.
Let its proposal be $q_d(x)=\prod_{j=1}^d q_1(x_j)$, with $p_1\ll q_1$.
Independence yields
\begin{align}
 1+\chi^2(p_d\Vert q_d)
 &=\int\frac{p_d(x)^2}{q_d(x)}\,dx
   =\prod_{j=1}^d\int\frac{p_1(x_j)^2}{q_1(x_j)}\,dx_j \\
 &=\{1+\chi^2(p_1\Vert q_1)\}^d.
 \label{eq:bf-is-smc-product-chi}
\end{align}
This is exponential growth in the second moment of the importance ratio
for any fixed nonzero coordinate-wise mismatch. By contrast, KL divergence
is additive for product distributions. Statements about exponential weight
variance should not be transferred to KL without a separate argument.

Annealing distributes this discrepancy over intermediate steps.
For a smooth path obtained by tempering the likelihood, let $\ell=\log L$ and define
$A(\beta)=\log Z_\beta$. Wherever differentiation under the integral is
valid, $A'(\beta)=\E_{p_\beta}\ell$ and
$A''(\beta)=\Var_{p_\beta}(\ell)$. A direct ratio calculation gives
\begin{align}
 1+\chi^2(p_{\beta+\delta}\Vert p_\beta)
 &=\frac{Z_{\beta+2\delta}Z_\beta}{Z_{\beta+\delta}^2},\\
 \log\{1+\chi^2(p_{\beta+\delta}\Vert p_\beta)\}
 &=\delta^2\Var_{p_\beta}(\ell)+O(\delta^3).
 \label{eq:bf-is-smc-local-chi}
\end{align}
The expansion requires the relevant nearby normalizers and moments,
including at $\beta+2\delta$, to exist. It suggests smaller increments
where the log likelihood varies more. If that variance grows proportionally
to $d$, the local calculation suggests increments of order $d^{-1/2}$.
It says nothing by itself about the cost of drawing an adequate sample
from each intermediate distribution.

\citet{beskos2014highdimsmc} prove stability results for product targets
with coordinate-wise kernels and appropriately increasing numbers of
intermediate distributions under drift and regularity assumptions. Their
discussion distinguishes those proofs from conjectured extensions to
full-dimensional proposal kernels. The more recent preprint of
\citet[Sections 2--3]{lefay2026smccomplexity} gives finite-error complexity
results for standard and waste-free SMC under controlled adjacent
chi-square discrepancy and explicit mixing assumptions. Its bounds retain
dependence on mixing times or spectral gaps; specialized results also
require conditions such as log-concavity and smoothness.

These results explain how SMC can avoid the direct importance-sampling
collapse in favorable settings. They do not establish a dimension-free
algorithm for arbitrary multimodal distributions. If an intermediate
kernel's spectral gap becomes exponentially small, treating that gap as a
constant would conceal the principal computational difficulty. Narrow
connections between modes, abrupt changes in relative mode mass, and
unrepresented basins therefore remain empirical questions even when
incremental ESS stays high \citep[Sections 3.1--3.3]{dai2022smc}.

The q20 example separates two dimensions that are easily confused. Its
parameter posterior has four free coordinates. Its recurrent state has
60 coordinates, of which 20 receive fresh process innovations, and each
observation is scalar. Annealed SMC over the UKF parameter posterior faces
the first geometry. An inner particle filter for the original state model
faces the second, including dependence on its history. Neither the label
``q20'' nor the total state dimension alone predicts their relative costs.

\section{State and parameter inference with SMC\texorpdfstring{$^2$}{2}}
\label{sec:bf-is-smc2}

The likelihood used by the current parameter sampler is a deterministic
UKF approximation. A different objective is Bayesian inference for the
specified generative state model itself. Let $\theta$ have prior $r$,
let $S_0$ have conditional distribution $\mu_\theta$, and let
$F_\theta(ds_t\mid s_{t-1})$ be the transition kernel. Observations have
conditional density $g_\theta(y_t\mid s_t)$. The joint posterior measure is
\begin{equation}
 p(d\theta,ds_{0:T}\mid y_{1:T})\propto
 r(\theta)\,d\theta\,\mu_\theta(ds_0)
 \prod_{t=1}^T F_\theta(ds_t\mid s_{t-1})
                         g_\theta(y_t\mid s_t).
 \label{eq:bf-is-smc2-joint}
\end{equation}
Writing the state transition as a measure accommodates deterministic state
coordinates; a nonsingular density on every state coordinate is not assumed.
Integrating the state path defines $L_T(\theta)=p(y_{1:T}\mid\theta)$.
Replacing this integral by a UKF likelihood changes the posterior.
Increasing a Monte Carlo particle count in an eligible particle filter
instead controls sampling error for the original integral.

\subsection{The inner filter and its likelihood estimate}

For clarity, consider a bootstrap particle filter with $N_x$ state particles
and resampling at every observation. At date $t$, draw predictive particles
from the transition after resampling the preceding filtered population;
weight them by $g_\theta(y_t\mid S_t^n)$. The likelihood estimate is
\begin{equation}
 \widehat L_t(\theta,U)
 =\prod_{s=1}^t\left\{\frac{1}{N_x}
                     \sum_{n=1}^{N_x}g_\theta(y_s\mid S_s^n)\right\},
 \qquad \E_{m_{\theta,N_x}}\widehat L_t(\theta,U)=L_t(\theta).
 \label{eq:bf-is-smc2-likelihood-estimate}
\end{equation}
Here $U$ contains the particle and ancestry variables and
$m_{\theta,N_x}$ is their simulation law. Unbiasedness follows from the
unnormalized particle-measure recursion and conditionally unbiased
resampling, not from independence of the factors in the product.
For example, let $\widehat\gamma_{s-1}(f)
=\widehat L_{s-1}\sum_n W_{s-1}^n f(S_{s-1}^n)$ and define
$Q_s f(x)=\int F_\theta(ds\mid x)g_\theta(y_s\mid s)f(s)$.
Conditional on the preceding particle history $\mathcal F_{s-1}$,
unbiased resampling and propagation give
\begin{equation}
 \E[\widehat\gamma_s(f)\mid\mathcal F_{s-1}]
 =\widehat L_{s-1}\sum_n W_{s-1}^n Q_sf(S_{s-1}^n)
 =\widehat\gamma_{s-1}(Q_sf).
 \label{eq:bf-is-smc2-unbiased-recursion}
\end{equation}
Starting from unbiased initial particles and iterating expectations yields
$\E\widehat\gamma_s(f)=\gamma_s(f)$. Taking $f=1$ proves likelihood
unbiasedness even though successive normalizer estimates are dependent.
Other proposals and adaptive resampling need their corresponding importance
factors and normalizer updates; the simple arithmetic averages above do
not apply unchanged. The logarithm of $\widehat L_t$ is generally biased
for $\log L_t$.

SMC$^2$ represents the parameter posterior with $N_\theta$ outer particles,
each carrying its own inner particle filter. As each observation arrives,
the inner filter advances and its likelihood increment updates the outer
weight. Resampling alone would repeatedly duplicate static parameter
values. Particle-MCMC mutation supplies new parameter values while
preserving an appropriate extended target
\citep[Sections 3.1--3.2]{chopin2013smc2}.

For a fixed inner particle count, that target is
\begin{equation}
 \widetilde\pi_{t,N_x}(\theta,u)
 =\frac{r(\theta)m_{\theta,N_x}(u)
                    \widehat L_t(\theta,u)}{p(y_{1:t})}.
 \label{eq:bf-is-smc2-extended}
\end{equation}
Integrating $u$ and using \eqref{eq:bf-is-smc2-likelihood-estimate}
gives $p(\theta\mid y_{1:t})$. This is the central identity in
\citet[Equation 7 and Proposition 1]{chopin2013smc2}. It is an identity
for the target measure at any eligible $N_x$, not a statement that a
finite outer particle population equals the posterior exactly.

For particle marginal Metropolis--Hastings (PMMH), propose
$\theta'\sim a(\cdot\mid\theta)$ and independently generate an inner
filter $u'\sim m_{\theta',N_x}$. The simulation laws cancel in the
Metropolis ratio, leaving
\begin{equation}
 \alpha((\theta,u),(\theta',u'))=1\wedge
 \frac{r(\theta')\widehat L_t(\theta',u')a(\theta\mid\theta')}
      {r(\theta)\widehat L_t(\theta,u)a(\theta'\mid\theta)}.
 \label{eq:bf-is-smc2-pmmh}
\end{equation}
On rejection, both the current parameter and its current likelihood estimate
are retained. Recomputing the denominator from an independent filter on every
attempt is a different algorithm. No unbiased log likelihood or target
derivative is needed for this PMMH construction.

\subsection{Adapt the inner accuracy to the mutation problem}

Inner particle count controls a tradeoff. A noisy likelihood estimate can
give an accepted state an unusually high estimated likelihood, making it
difficult to leave. Increasing $N_x$ reduces that noise in favorable
regimes, but multiplies the cost of each outer proposal. A parameter move
at date $t$ normally reruns a state filter through the complete prefix,
costing order $tN_x$ particle operations per proposal. Parallel execution
across parameters reduces elapsed time without removing this total work.

The original SMC$^2$ paper permits increases of $N_x$ and studies complexity
under restrictive mixing and posterior-concentration assumptions
\citep[Sections 3.6--3.7 and Appendix B]{chopin2013smc2}.
\citet[Sections 3--4]{botha2023adaptivesmc2} separate adaptation into three
decisions: when to change the inner population, how many particles to use,
and how to replace the stored particle system. Their movement-based
trigger avoids interpreting high acceptance from tiny proposals as
successful rejuvenation. They jointly adapt the number of PMMH iterations
and $N_x$, and compare likelihood-variance rules with a cost criterion
based on expected squared jumping distance (ESJD).

If a pilot supports the approximation
$\Var(\log\widehat L_t)\simeq v_t/N_x$, a measured variance $s^2$ at
count $N_x$ suggests
\begin{equation}
 N_x'\simeq N_x\frac{s^2}{v_{\rm target}}.
 \label{eq:bf-is-smc2-variance-allocation}
\end{equation}
This is a pilot allocation model, not an identity. It can fail when small
counts give highly skewed estimates or variance changes sharply across
parameter regions. The paper examines moderated updates and candidate
counts because direct rescaling can overshoot. Its movement-based method
estimates the number $R(N_x)$ of mutations needed to reach a declared
movement target and compares costs proportional to $N_xR(N_x)$
\citep[Section 4.2 and Algorithm 5]{botha2023adaptivesmc2}.
Neither a universal log-variance target nor the current population mean
is automatically adequate for a multimodal application: the mean may lie
between modes, and the likelihood noise may differ substantially among them.

\subsection{Changing the inner particle system}
\label{sec:bf-is-smc2-refresh}

Changing $N_x$ changes the auxiliary particle-system space. It must preserve
the intended parameter marginal and the correct new extended distribution.
At fixed $\theta$, write the old and new PF laws as $m_N(u\mid\theta)$
and $m_M(v\mid\theta)$. Proposing a fresh $v\sim m_M$ and using the
reverse simulation law $m_N$ gives the exchange importance factor
\begin{equation}
 \frac{r(\theta)m_M(v\mid\theta)\widehat L_M(\theta,v)
                         m_N(u\mid\theta)}
      {r(\theta)m_N(u\mid\theta)\widehat L_N(\theta,u)
                         m_M(v\mid\theta)}
 =\frac{\widehat L_M(\theta,v)}{\widehat L_N(\theta,u)}.
 \label{eq:bf-is-smc2-exchange}
\end{equation}
This correction can increase weight variability, but that cost does not
make it optional. A correctly constructed conditional particle-filter
update provides another route, retaining a selected trajectory and
changing the remaining auxiliary variables under the appropriate law
\citep[Section 3.6]{chopin2013smc2}.

The \texttt{replace} option in
\citet[Section 4.3]{botha2023adaptivesmc2} uses a different approximation.
The optimal backward kernel would involve the exact likelihood in its
normalizer. Replacing that likelihood by the new PF estimate makes the
formal incremental weight cancel to one. The authors explicitly describe
this as an approximation to the backward kernel. Its consequence can be
checked directly in the present notation: the proposed expression
\begin{equation}
 B(v,u)=\frac{m_N(u\mid\theta)\widehat L_N(\theta,u)}
                    {\widehat L_M(\theta,v)}
 \quad\hbox{satisfies}\quad
 \int B(v,u)\,du=\frac{L(\theta)}{\widehat L_M(\theta,v)},
 \label{eq:bf-is-smc2-replace-normalization}
\end{equation}
which is generally not one. It is therefore not an exactly normalized
backward Markov kernel at finite $M$.

The same issue appears without backward-kernel notation. Under the new
extended target, conditional on $\theta$, the auxiliary law must be
$m_M(v\mid\theta)\widehat L_M(\theta,v)/L(\theta)$.
An independent unweighted refresh draws $m_M(v\mid\theta)$ instead.
For an elementary example, let $L=1$ and let the simulation law give
$\widehat L=1/2$ or $3/2$ with equal probabilities. The extended law gives
those outcomes probabilities $1/4$ and $3/4$, while an unweighted refresh
returns $1/2$ and $1/2$. Subsequent invariant PMMH can reduce the discrepancy
through mixing, but does not prove exactness after an arbitrary finite
number of steps.

This is a limitation of the replacement operation, not a refutation of
adaptive particle allocation. The ESJD and variance ideas can be studied
with exact exchange or an appropriate conditional-filter update. Such a
combination must be described as an explicit variant; faithfully copying
an approximate operation does not make it exact.

\subsection{Switching parameter-update kernels}

PMMH proposes parameters while integrating state uncertainty through a
particle likelihood estimate. Particle Gibbs instead uses a selected
state trajectory as part of its conditioning structure and refreshes it
with conditional SMC. One can outperform the other in some stages and
lose that advantage later: PMMH can suffer from noisy likelihood ratios,
while Gibbs updates can mix slowly when parameters and trajectories are
strongly dependent.

\citet[Sections 3--4]{botha2023switchingsmc2} investigate switching between
these kernels, allowing different inner particle counts. Their criterion
uses movement in the least-moving parameter coordinate relative to
inner-particle cost, then allocates the remaining mutation effort to the
selected kernel. This addresses a weakness of fixing one rejuvenation
method for the whole sequence. The preprint's Appendix A, however,
substitutes particle likelihood estimates into some backward-kernel
normalizers when switching to PMMH. Those operations need the same
exactness distinction as Section~\ref{sec:bf-is-smc2-refresh}; the existence
of a valid PMMH kernel after a switch does not establish that its initial
auxiliary variables have the required distribution.

For the recurrent q20 model there is an additional structural difficulty.
Hidden and cell coordinates evolve deterministically given the preceding
state. A generic backward-sampling formula that evaluates a nonsingular
full-state transition density is therefore not directly applicable.
Possible approaches include conditioning on the deterministic recurrence
or expressing the model through its genuine innovations and initial state.
Either requires a derivation of the resulting conditional sampler. Adding
small artificial noise would change the state model and is not a harmless
way to make generic particle Gibbs executable.

\subsection{Two kinds of density tempering}

Annealed SMC and SMC$^2$ are not mutually exclusive. In addition to adding
observations sequentially, one can temper a full-data likelihood. But three
intermediate targets must be distinguished. The intended marginal path is
$r(\theta)L_T(\theta)^\beta$. The extended path
$r(\theta)m_{\theta,N_x}(u)\widehat L_T(\theta,u)^\beta$ has marginal
$r(\theta)\E_m[\widehat L_T^\beta]$. In general,
\begin{equation}
 \E_m[\widehat L_T^\beta]\ne L_T^\beta,
 \qquad 0<\beta<1,
 \label{eq:bf-is-smc2-power-inequality}
\end{equation}
with the left side no greater than the right by concavity when the
expectations exist. Properly corrected extended-space tempering can still
have the desired endpoint at $\beta=1$; it must not claim the marginal
intermediate distributions of ordinary likelihood annealing
\citep[Section 5.2]{chopin2013smc2}.

A third path tempers the conditional observation likelihood before
integrating out the states:
\begin{equation}
 \gamma_\beta^{\rm joint}(\theta,ds_{0:T})
 =r(\theta)\mu_\theta(ds_0)
     \prod_{t=1}^T F_\theta(ds_t\mid s_{t-1})
                         g_\theta(y_t\mid s_t)^\beta.
 \label{eq:bf-is-smc2-joint-tempering}
\end{equation}
Its parameter marginal involves $\E_{S\mid\theta}
[\prod_t g_\theta(y_t\mid S_t)^\beta]$, rather than
$\{\E_{S\mid\theta}\prod_t g_\theta(y_t\mid S_t)\}^\beta$.
A PF using $g^\beta$ estimates the former normalizer. This difference
underlies the compatible and incompatible kernel combinations discussed
in \citet[Section 2.3]{botha2023switchingsmc2}.

For a scalar Gaussian observation with variance $R$ and $\beta>0$,
\begin{equation}
 \calN(y;m,R)^\beta
 =(2\pi R)^{(1-\beta)/2}\beta^{-1/2}
                         \calN(y;m,R/\beta).
 \label{eq:bf-is-smc2-gaussian-power}
\end{equation}
Using observation variance $R/\beta$ without the prefactor computes a
different potential. If $R$ depends on $\theta$, the omission also changes
relative parameter weights. At $\beta=0$, use the limiting potential
$g^0=1$ directly rather than an infinite-variance Gaussian representation.

\subsection{An observation-adapted state proposal}
\label{sec:bf-is-smc2-adapted-q20}

In the preserved q20 model, conditional on the preceding full state
$s_{t-1}$, the hidden and cell updates are fixed. The stochastic coordinates
satisfy $z_t\sim\calN(m_\theta(s_{t-1}),Q)$ and
$y_t\mid z_t\sim\calN(H_\theta z_t+b_\theta,R)$, where $R>0$ is scalar
and $Q$ is positive definite. The joint covariance of $(z_t,y_t)$ is
\begin{equation}
 \begin{pmatrix}Q&QH_\theta^\top\\H_\theta Q&S_\theta\end{pmatrix},
 \qquad S_\theta=H_\theta QH_\theta^\top+R.
 \label{eq:bf-is-smc2-adapted-joint-cov}
\end{equation}
Conditioning this Gaussian gives $K_\theta=QH_\theta^\top S_\theta^{-1}$
and
\begin{equation}
 z_t\mid s_{t-1},y_t,\theta\sim
 \calN\!\left(m+K_\theta(y_t-H_\theta m-b_\theta),
             Q-QH_\theta^\top S_\theta^{-1}H_\theta Q\right).
 \label{eq:bf-is-smc2-adapted-proposal}
\end{equation}
The predictive observation density for ancestor $i$ is
$a_i=\calN(y_t;H_\theta m_\theta(s_{t-1}^i)+b_\theta,S_\theta)$.
With preceding normalized weights $w_i$, a fully adapted auxiliary filter
selects that ancestor with probability $w_i a_i/\sum_jw_j a_j$ and draws
from \eqref{eq:bf-is-smc2-adapted-proposal}. Bayes' formula makes its
post-proposal importance ratio constant:
\begin{equation}
 \frac{w_i\,p(z_t\mid s_{t-1}^i,\theta)
                    g_\theta(y_t\mid z_t)}
      {[w_i a_i/\sum_jw_j a_j]\,
                    p(z_t\mid s_{t-1}^i,y_t,\theta)}
 =\sum_jw_j a_j.
 \label{eq:bf-is-smc2-adapted-weight}
\end{equation}
That weighted predictive average is the likelihood-normalizer increment.
This derivation by Gaussian conditioning exploits the actual stochastic
coordinates and preserves the deterministic recurrence. Its potential
benefit is avoiding proposals inconsistent with the current observation.
It does not prevent loss of ancestral histories or guarantee efficient
parameter-mode exploration. No implementation or performance result for
this adapted q20 filter is asserted here.
